\documentclass[12pt]{article}
\usepackage[utf8]{inputenc}
\usepackage[T1]{fontenc}
\usepackage[margin=1in]{geometry}
\usepackage{setspace}
\usepackage{xcolor}
\usepackage[colorlinks=true, urlcolor=blue!60!black, linkcolor=blue!60!black]{hyperref}
\setlength{\parindent}{2em}
\setlength{\parskip}{1.25em}
\renewcommand{\baselinestretch}{1.0}
\begin{document}
\begin{center}
{\LARGE \textit{Clapper v.\ Amnesty International USA}, 568 U.S.\ 398 (2013)\par}
\vspace{0.8em}
{\large Case Brief --- Four-Page Edition\par}
\end{center}
\vspace{1em}

\section{Memory Jogger}

Lawyers, journalists, and human rights groups lacked standing to challenge Section 702 of FISA (50 U.S.C.\ \S~1881a): their fear of interception rested on a five-link ``speculative chain of possibilities,'' not ``certainly impending'' harm, and they could not ``manufacture standing'' through costly precautions (\href{https://supreme.justia.com/cases/federal/us/568/398/}{Justia}).

\section{Detailed Case Facts}

The 1978 Foreign Intelligence Surveillance Act required the FISC to find probable cause that a target was a foreign power or its agent and that the monitored facilities were used by that target. After the post-9/11 warrantless surveillance program, Congress added Section 702 in 2008. It lets the Attorney General and the Director of National Intelligence authorize, for up to a year, targeting of non-U.S.\ persons reasonably believed to be abroad, without individualized probable cause or identified facilities, subject to FISC approval of targeting and minimization procedures; it bars intentionally targeting U.S.\ persons. Respondents were U.S.\ attorneys and human rights, labor, legal, and media organizations who communicated with foreign clients and sources, including attorney Scott McKay (whose client's prior communications had been intercepted some 30,000 times), attorney Sylvia Royce, journalist Christopher Hedges, and researcher Joanne Mariner. They alleged that the law forced them to stop some communications, speak in generalities, and travel abroad to meet in person.

\section{Procedural History}

Respondents sued the day the law was enacted. The district court granted summary judgment for the government for lack of standing. The Second Circuit reversed, finding an ``objectively reasonable likelihood'' of interception and present protective costs, and denied rehearing en banc 6--6. The Supreme Court reversed on February 26, 2013.

\section{Judicial Votes}

Reversed and remanded, \textbf{5--4}.
\begin{itemize}
\item \textbf{Majority:} Justice Alito, joined by Chief Justice Roberts and Justices Scalia, Kennedy, and Thomas.
\item \textbf{Concurrences:} None.
\item \textbf{Dissent:} Justice Breyer, joined by Justices Ginsburg, Sotomayor, and Kagan.
\end{itemize}

\section{Holding}

A threatened injury must be ``certainly impending'' (or, per footnote 5, pose a ``substantial risk''), not merely an ``objectively reasonable likelihood'' resting on an attenuated chain involving independent decisionmakers. Plaintiffs cannot manufacture standing through self-inflicted costs responding to non-imminent harm, and such costs are not fairly traceable to the statute. Footnote 5's ``substantial risk'' language (later elevated in \textit{Susan B.\ Anthony List v.\ Driehaus}) and the Court's assurances about other avenues of review are dicta.

\section{Analysis of Opinions}

\textbf{Majority (Alito).} Standing review is ``especially rigorous'' when it would require reviewing the political branches' intelligence activities. Respondents' theory required that (1) the government target their foreign contacts; (2) it use Section 702 rather than other authorities; (3) FISC judges approve; (4) interception succeed; and (5) respondents be party to the intercepted communications. At summary judgment, \textit{Lujan} requires specific facts, and the second link also defeated traceability. Protective costs cannot substitute: otherwise any plaintiff could, ``for the price of a plane ticket,'' reduce Article III to a ``nonparanoid fear.'' Because McKay's client was already monitored under older FISA, precautions were not traceable to Section 702, and a ``subjective `chill' '' is insufficient under \textit{Laird v.\ Tatum}. \textit{Laidlaw}, \textit{Meese v.\ Keene}, and \textit{Monsanto} were distinguishable. The lack of any other plaintiff does not create standing (\textit{Valley Forge}), and criminal-defendant notice and provider challenges preserve review.

\textbf{Dissent (Breyer).} Interception was as likely as most future events that common sense predicts. Respondents' communications contained foreign intelligence information that Section 702 newly allowed the government to collect; both sides had strong motives; the government's past monitoring showed its appetite; and the NSA's capacity and the FISC's 99.9\% approval rate made interception virtually assured. ``[C]ertainty is not, and never has been, the touchstone of standing''; precedent and traditional anticipatory actions require only a reasonable or high probability, and mitigation costs independently suffice.

\section{Examples: Future Applications}

\begin{enumerate}
\item \textbf{Same-Side: Ransomware Without Exfiltration.} A patient buys credit monitoring after an attack that took no data. Her fear rests on speculative third-party choices, and the purchase cannot manufacture standing.
\item \textbf{Same-Side: Facial Recognition.} A journalist who avoids protests after a law authorizes police facial-recognition software alleges only a subjective chill.
\item \textbf{Opposite-Side: Pre-Enforcement Speech Threat.} A group previously hauled before a commission, and already threatened with a new complaint, faces a ``substantial risk'' under \textit{Driehaus}.
\item \textbf{Opposite-Side: Data Stolen and Misused.} After hackers steal Social Security numbers and start filing fraudulent returns, the risk is substantial and mitigation costs are traceable.
\item \textbf{Fence-Sitter: Upstream Collection.} A news organization knows its traffic crosses tapped cables but cannot prove retention of any message; compare \textit{Wikimedia Foundation v.\ NSA} (4th Cir.\ 2021).
\end{enumerate}

\section{Critique}

\textbf{Article III.} Supporters say rigorous standing confines courts to real ``Cases'' and ``Controversies.'' Some originalist and textualist scholars, however, see injury in fact as an ahistorical twentieth-century gloss that conflates jurisdiction with the merits.

\textbf{Catch-22.} Critics say plaintiffs cannot prove secret surveillance without discovery and cannot get discovery without proof. The Court's notice assurance proved inaccurate: prosecutors had been withholding notice until the Justice Department issued its first FAA notice in October 2013. Critics also contrast \textit{Driehaus}'s laxer treatment of probabilistic harm.

\section{Key Quotations}

\begin{quote}
``In other words, respondents cannot manufacture standing merely by inflicting harm on themselves based on their fears of hypothetical future harm that is not certainly impending.''
\end{quote}
\noindent --- Justice Alito, for the Court, 568 U.S.\ at 416.

\begin{quote}
``We decline to abandon our usual reluctance to endorse standing theories that rest on speculation about the decisions of independent actors.''
\end{quote}
\noindent --- Justice Alito, for the Court, 568 U.S.\ at 414.

\begin{quote}
``Our cases do not uniformly require plaintiffs to demonstrate that it is literally certain that the harms they identify will come about. In some instances, we have found standing based on a `substantial risk' that the harm will occur, which may prompt plaintiffs to reasonably incur costs to mitigate or avoid that harm.''
\end{quote}
\noindent --- Justice Alito, for the Court, 568 U.S.\ at 414 n.5.

\end{document}
